%% ============================================================================
\section{Probabilistic Logical Coherence}
\label{sec:framework}
%% ============================================================================

In this section, we formalize the notion of logical coherence, define a probabilistic model for it, 
and introduce the Logical Coherence Score (LCS). Table~\ref{tab:notation} (App.~\ref{app:framework})
lists the notation. 

\subsection{Claims and the Committed World}
\label{sec:setup}

A response is decomposed into atomic claims $\A=(a_1,\dots,a_n)$; a \emph{world}
$x\in\{0,1\}^n$ assigns each claim a truth value, with $x_i=1$ read as ``$a_i$ holds''. A claim is a
statement the response mentions, whether it asserts it or denies it.

%: both responses of Figure~\ref{fig:voyager} mention $a_4$ in order to reject it.

\begin{definition}[Committed world]
\label{def:committed}
The \emph{committed world} $\xs\in\{0,1\}^n$ assigns each claim the truth value under which the
argument is evaluated. Under \emph{verdict polarity} (our default), $\xs_i$ is a factuality
pipeline's verdict on claim $i$: $1$ if supported, $0$ if not. Under \emph{assertion polarity},
$\xs_i$ is the response's own stance: $1$ if it asserts claim $i$, $0$ if it denies it. We call the
two the \emph{verdict world} and the \emph{asserted world}, respectively.
\end{definition}

Any decompose-then-verify pipeline can supply verdict polarity \citep{factscore,safe,veriscore}, or
FactReasoner's posteriors \citep{factreasoner} thresholded at $\frac12$, so responses over the same
claims are judged against one world. The polarities coincide when a response asserts exactly the
supported claims, as both Voyager responses do ($\xs=(1,1,1,0,0)$).

We further assume that every claim $a_i$ has a prior probability $\pi_i\in(0,1)$, with
$\pi(x)=\prod_i\pi_i^{x_i}(1-\pi_i)^{1-x_i}$. Label-only factuality pipelines give flat or uniform priors,
$\pi_i=\frac12$, which we use throughout. A probabilistic factuality pipeline's posteriors can serve as priors,
provided they are computed without the response's own relations. Every result below holds for any
priors (App.~\ref{app:twostage}).

\subsection{Relations and Their Sign}
\label{sec:taxonomy}

While discourse relations describe how claims are connected in text, a probabilistic model requires 
their logical semantics. We represent a relation by the joint truth assignments of its claims that 
it excludes. We focus on binary relations only.

\begin{definition}[Coupling and relation]
\label{def:coupling}
\label{def:relation}
A \emph{coupling} $\kappa$ is identified with the set $V(\kappa)\subseteq\{0,1\}^2$ of joint
truth values $(x_s,x_t)$ of its endpoints $a_s,a_t$ that \emph{violate} it: entailment \ent{}
($a_s\Rightarrow a_t$, $V=\{(1,0)\}$), equivalence \eqv{} ($\{(0,1),(1,0)\}$), co-necessity \cnc{}
(at least one holds, $\{(0,0)\}$), contradiction \con{} (not both, $\{(1,1)\}$), exclusion \exc{}
(exactly one, $\{(0,0),(1,1)\}$), and \emph{none} ($\varnothing$). A \emph{relation}
$r=(e_r,\kappa_r,p_r)$ has a scope $e_r=(s_r,t_r)$ of two distinct claims, a coupling $\kappa$, and a weight
$p_r\in[0,1)$. A \emph{relation graph} $\R$ is a finite multiset of relations. We write
$x_e=(x_s,x_t)$ for the values (or projection) of a world on a scope.
\end{definition}

% Keep it commented out for rebuttal (reviewers may ask for it, but it is not needed in the main text).
% Relation weights stay below $1$ so that no relation becomes a hard constraint and every world
% keeps positive probability in the model of \S\ref{sec:mrf}. At $p_r=1$, a single conflict would
% drive the score to its minimum whatever the rest of the graph, and strict monotonicity would fail
% (App.~\ref{app:weights}).

Relations are annotated with a discourse \emph{sense} from PDTB~3.0 and RST \citep{pdtb,mann},
which a deterministic \emph{compile map} associates with a coupling (Table~\ref{tab:couplings},
App.~\ref{app:framework}). Only
\sense{Refutation}, an asserted incompatibility, compiles to \con{}. Furthermore, \sense{Contrast} and
\sense{Concession}, whose arguments are typically both true (``sales rose, but profits fell''),
compile to \emph{none}, as do the temporal senses.

\begin{definition}[Support and conflict]
\label{def:sign}
A relation $r$ with $V(\kappa_r)\neq\varnothing$ is a \emph{support} if the committed world
satisfies it, $\xs_{e_r}\notin V(\kappa_r)$, and a \emph{conflict} if it violates it. They form the
multisets $\Rsup$ and $\Rcon$, with $\R=\Rsup\uplus\Rcon$. Relations with $V=\varnothing$ produce no
factor and are dropped.
\end{definition}

A contradiction between two truths is a conflict, but between a truth and a falsehood a support,
since it refutes the falsehood; an entailment from a truth to a falsehood is a conflict. Response~B in
Figure~\ref{fig:voyager} has $\Rsup=\{a_1\Rightarrow a_2,\
a_3\con a_5\}$ and $\Rcon=\{a_2\con a_3,\ a_1\Rightarrow a_4\}$, while all of Response~A's relations are supports.
The sign depends only on the verdicts, so an inference between two false claims counts as a
support: whether premises are true is the factuality score's job.

\subsection{Estimating Relations}
\label{sec:mining}

In practice, the relations are estimated from the response. A relation's weight is the product of a
\emph{confidence} $c_r$ that the response draws it and the \emph{strength} $g_r$ with which it does,
\begin{equation}
p_r=c_r\cdot g_r,\qquad c_r,g_r\in[0,1],
\label{eq:decomp}
\end{equation}
clipped to at most $1-10^{-6}$. A \emph{pairwise} miner scores claim pairs within a local window, with $c_r$ from label-token
probabilities and $g_r$ from a \texttt{Yes}/\texttt{No} judgment \citep{kadavath}; a \emph{graph-level}
miner extracts all relations, with sense, direction and strength, in one call. 
Over $k=5$ samples it keeps relations found by a strict majority, with $c_r$ the vote share and $g_r$
from the strength band. Both miners recover only relations the response asserts
(App.~\ref{app:estimation}).



\subsection{The Coherence Markov Random Field}
\label{sec:mrf}

The coherence Markov random field (MRF) has one binary variable per claim, whose value $x_i=1$
means that claim $a_i$ is true. Each claim is assigned a prior factor $\phi_i(x_i)=\pi_i^{x_i}(1-\pi_i)^{1-x_i}$
and each relation $r$ contributes a factor $\psi_r$ that penalizes exactly the joint truth assignments that the relation forbids:
\begin{equation}
\psi_r(u,v)=
\begin{cases}
1-p_r & \text{if } (u,v)\in V(\kappa_r),\\
1 & \text{otherwise},
\end{cases}
\label{eq:psi}
\end{equation}

The resulting distribution over worlds is:
\begin{equation}
\PR(x)=\frac{\pi(x)\prod_{r\in\R}\psi_r(x_{e_r})}{Z_\R},
\label{eq:mrf}
\end{equation}
where $Z_\R=\sum_{x\in\{0,1\}^n}\pi(x)\prod_{r\in\R}\psi_r(x_{e_r})$ is the partition function
that normalizes the distribution, a sum over all $2^n$ worlds that equals the prior expectation
$\E_\pi\bigl[\prod_{r\in\R}\psi_r\bigr]$.
Each relation leaves a world that satisfies it unchanged and multiplies the probability of one that
violates it by $1-p_r$.

% Keep it commented out for rebuttal (reviewers may ask for it, but it is not needed in the main text).
% This construction also gives the key property needed for our monotonicity results.
% Since each relation factor takes only two values, $1$ and $1-p_r$, the committed-world 
% assignment is always either its most preferred or least preferred configuration. 
% In the terminology of App.~\ref{app:why}, every relation factor is therefore \emph{sign-definite}, 
% which is precisely the condition required for coherence-improving edits to increase the 
% score monotonically.


\subsection{The Coherence Order and the Score}
\label{sec:order}
\label{sec:score}

\begin{definition}[Coherence order]
\label{def:order}
An \emph{elementary improvement} of $\R$ adds a support, raises a support's weight, deletes a
conflict, or lowers a conflict's weight. We say that $\R'\coh\R$ if $\R'$ is reached from $\R$ by finitely many
elementary improvements, and $\R'\scoh\R$ if at least one of them changes a weight by a non-zero
amount, counting an addition or a deletion as a change from or to weight $0$.
\end{definition}

% Edits that improve one part of an argument and degrade another stay unranked. 

For a fixed $\xs$, we require a score $S(\R;\pi)$ to be \textbf{(A1) monotone},
$\R'\scoh\R\Rightarrow S(\R')>S(\R)$ for every $\pi$; \textbf{(A2) MRF-grounded}, depending on $\R$ only through
$\PR$ and $\pi$; and \textbf{(A3) neutral}, with $S(\varnothing;\pi)$ independent of $\pi$. A1
and A3 alone admit simpler scores such as a signed sum of weights; A2, a modeling choice, makes
relations interact through $Z_\R$ (App.~\ref{app:why}).

% A1 is demanding because the edited relation can sit anywhere in an
% arbitrary graph. It suggests reading the MRF at the world the argument is about, before and after
% the relations.

\begin{definition}[Logical coherence score]
\label{def:score}
\begin{equation}
\LCS(\R)=\sigma\bigl(\operatorname{logit}\PR(\xs)-\operatorname{logit}\pi(\xs)\bigr),
\label{eq:lcs}
\end{equation}
with $\sigma$ the logistic function and $\operatorname{logit}u=\log\frac{u}{1-u}$.
\end{definition}

The log-odds difference in \eqref{eq:lcs} is Good's \emph{weight of evidence} \citep{good1950} that 
the relations contribute in favor of the complete truth assignment $\xs$. $\LCS(\R)$ is the probability of $\xs$ for a reader who starts from even odds and updates on the
relations: $\frac12$ when they neither confirm nor disconfirm it (as with no relations), tending to $0$
when they rule out $\xs$ and to $1$ when they force it. It measures an amount of evidence, not an
average.


The results are easier to understand in terms of:
\begin{equation}
\Lambda=\log\frac{\PR(\xs)}{\pi(\xs)}=\sum_{r\in\Rcon}\log(1-p_r)-\log Z_\R,
\label{eq:lambda}
\end{equation}
since $\psi_r(\xs_{e_r})$ is $1$ for a support and $1-p_r$ for a conflict. For a fixed claim set and
prior, $\LCS$ is strictly increasing in $\Lambda$ (App.~\ref{app:lcs-lambda}), and both need $Z_\R$, which
is \#P-hard in general \cite{koller2009probabilistic}.



\begin{lemma}[Edit identity]
\label{lem:edit}
For any $\R$, relation $r$ and priors,
$\Lambda(\R\uplus\{r\})-\Lambda(\R)=\log\psi_r(\xs_{e_r})-\log\E_{\PR}[\psi_r(x_{e_r})]$. Let
$q_r=\PR(x_{e_r}\in V(\kappa_r))\in(0,1)$; then this is $-\log(1-p_rq_r)>0$ for a support and
$\log\frac{1-p_r}{1-p_rq_r}<0$ for a conflict, strictly when $p_r>0$.
\end{lemma}

The lemma is the virtual-evidence update \citep{pearl1988,chan2005}, and the MRF a Markov logic
network \citep{mln} with clause weights $-\log(1-p_r)$; what is new is the reading at $\xs$. A relation's effect depends on the context in magnitude, through $q_r$, but not in sign, so the main
result follows.

\begin{theorem}[Monotonicity]
\label{thm:mono}
Fix $\xs$. For all priors $\pi\in(0,1)^n$ and all relation graphs $\R,\R'$ with weights in $[0,1)$:
if $\R'\coh\R$ then $\LCS(\R')\ge\LCS(\R)$, and if $\R'\scoh\R$ then $\LCS(\R')>\LCS(\R)$. The
same holds for $\Lambda$.
\end{theorem}

The proof does not use the fact that scopes are pairs. It ranks the Voyager pair (A is reached from B
by elementary improvements; $\LCS=\vyALCS$ against $\vyBLCS$), and the shared refutation $a_3\con a_5$
adds $\vyADFour$ nats to $\Lambda$ in A but only $\vyBDFour$ in B, where $a_2\con a_3$ has made $a_3$
less probable.

\subsection{Per-Relation Evidence and Interaction}
\label{sec:decomp}

$\Lambda$ separates into what each relation contributes alone and what the relations contribute
together. Let $v_r=\pi(x_{e_r}\in V(\kappa_r))$ be the prior mass of the assignments $r$ forbids
($\frac14$ for \ent, \con, \cnc{} and $\frac12$ for \eqv, \exc{} under flat priors).

\begin{proposition}[Decomposition]
\label{prop:decomp}
For any $\R$ and priors, $\Lambda(\R)=\sum_{r\in\R}\lambda_r-I_\R$, where $\lambda_r=\Lambda(\{r\})$
equals $-\log(1-p_rv_r)$ for a support and $\log\frac{1-p_r}{1-p_rv_r}$ for a conflict, and
$I_\R=\log Z_\R-\sum_r\log(1-p_rv_r)$. $I_\R$ is a sum over the connected components of the
relation graph, so it vanishes when no two relations share a claim.
\end{proposition}

$\lambda_r$ is Lemma~\ref{lem:edit} with no other relation present. $I_\R$ depends only on the priors
and the graph, not on $\xs$; for two relations its sign is that of $\pi(V_1\cap V_2)-v_1v_2$
(App.~\ref{app:proof-score}). When the forbidden assignments go together (a repeated relation;
$a\Rightarrow b$ beside $a\con w$) $I_\R>0$, so supports earn less, a repeated support saturating, and
conflicts cost more, a repeated violation compounding; when they exclude each other (a chain; the two
halves of an equivalence) the reverse holds. Since the verdicts do not enter it, refuting a
falsehood by citing $a$ also discounts the inferences drawn from $a$ (\S\ref{sec:limitations}).

\paragraph{A hierarchy of simpler scores.} For small weights the interaction vanishes and
$\lambda_r$ is $p_rv_r$ for a support and $-p_r(1-v_r)$ for a conflict, so $\Lambda=\Llin+O(p^2)$ with
\begin{equation}
\Llin=\sum_{r\in\Rsup}p_rv_r-\sum_{r\in\Rcon}p_r(1-v_r).
\label{eq:llin}
\end{equation}
Three $O(|\R|)$ scores satisfying A1 and A3 add the ingredients of $\Lambda$ one at a time: the
\emph{signed weight sum} $\Ssum=\sum_{\Rsup}p_r-\sum_{\Rcon}p_r$ ($2\Llin$ with every $v_r=\frac12$) has
none; the \emph{linear evidence} $\Llin$ adds the asymmetry (a violated \ent{} costs three times what a
satisfied one earns); and the \emph{independent evidence} $\Lambda^{\perp}=\sum_r\lambda_r$ adds the
curvature (the ratio is $\dcRatioOne$ at $p_r=\dcP$, the strong band), treating relations as separate
evidence as naive Bayes does. $\Llin$ scores a constraint split into halves of equal weight as the
whole; the curvature breaks this and the interaction restores it, and only $\Lambda$ satisfies A2. On
LoCoBench $\Llin$ is at least as accurate as $\LCS$ on every arm and on natural text the two cannot be
told apart (\S\ref{sec:experiments}), so we recommend $\Llin$ in practice; $\LCS$ is the model it
approximates, in which A2 holds and redundant and contested evidence are treated exactly
(Table~\ref{tab:three}; Figure~\ref{fig:decomp}, App.~\ref{app:decomp}).

% App.~\ref{app:props} shows that the result rests on two choices that cannot be changed,
% sign-definite tables and the committed world as the readout, and collects further properties of
% the score: additivity across independent sub-arguments, bounds, and saturation of redundant
% support.
% Sec. 2.6 (Why This Score) moved to App. app:props; restore from there if reviewers ask.
